% Generated draft: requires mathematical and source review.
\documentclass[11pt]{article}
\usepackage[margin=25mm]{geometry}
\usepackage{fontspec}
\setmainfont{Latin Modern Roman}
\usepackage{amsmath,amssymb,mathtools,unicode-math}
\setmathfont{Latin Modern Math}
\usepackage{enumitem,needspace,xurl,hyperref}
\setlength{\parindent}{0pt}
\setlength{\parskip}{5pt}
\setlength{\emergencystretch}{3em}
\newcommand{\N}{\mathbb{N}}
\newcommand{\Z}{\mathbb{Z}}
\newcommand{\Q}{\mathbb{Q}}
\newcommand{\R}{\mathbb{R}}
\newcommand{\C}{\mathbb{C}}
\newcommand{\F}{\mathbb{F}}
\newcommand{\A}{\mathbb{A}}
\newcommand{\PP}{\mathbb{P}}
\newcommand{\E}{\mathbb{E}}
\newcommand{\eps}{\varepsilon}
\newcommand{\dd}{\,\mathrm{d}}
\begin{document}
\section*{UnsolvedMath ID: 30005328}
\textbf{Source problem code:} catalogue entry
\subsection{Definitions and mathematical statement}
For a polynomial \(P(q)=\sum_{i\geq 0}a_iq^i\) with real coefficients, extend its coefficient sequence by zeros beyond its degree. Following the supplied definition, call \(P\) \emph{log-concave} if
\[
a_i^2\geq a_{i-1}a_{i+1}\qquad\text{for every integer }i\geq 1.
\]
Here log-concavity concerns the coefficients as the exponent of \(q\) varies.

The primary question, attributed in the extract to Conjecture 3 and to reference [3, Conjecture 7.4], is the following. Writing the source's \(q\)-Stirling polynomials as
\[
S_q[n,k]=\sum_{i\geq0}a_{n,k,i}q^i,
\]
is it true that, for every admissible pair \((n,k)\),
\[
a_{n,k,i}^{\,2}\geq a_{n,k,i-1}a_{n,k,i+1}
\qquad\text{for every integer }i\geq1?
\]
\textbf{Unresolved specification.} The supplied extract does not define this family \(S_q[n,k]\), give its recurrence or initial conditions, or explicitly state its index domain. Thus the displayed question preserves the source's target but does not uniquely determine the polynomials or the quantification over \((n,k)\). No convention for these missing data is imposed here.

The extract reports computer checks for all \(n,k\leq50\), without specifying the lower bounds or valid index pairs in that sentence. It also reports that, for \(k\geq2\), asymptotic normality of the coefficients as \(n\to\infty\) has been confirmed, citing [3, Theorem 7.7]. The probability distribution, normalization, and treatment of \(k\) in that limit are not supplied, so this report is retained without a reconstructed theorem statement. The extract records attempts using Lorentzian polynomials [7], says that the theory of atlases [6] had not been tried, and associates a separate question about \(q\)-log-concavity with ideas in [10]. These remarks do not specify an additional precise conjecture.

The extract also gives a related type \(B\) extension from set partitions to signed set partitions of type \(B_n\):
\[
S_q^B[n,k]=S_q^B[n-1,k-1]+[2k+1]_qS_q^B[n-1,k],
\qquad
S_q^B[0,k]=
\begin{cases}
1,&k=0,\\
0,&\text{otherwise}.
\end{cases}
\]
The conventional interpretation of the bracket is \([m]_q=\sum_{j=0}^{m-1}q^j\) for positive integers \(m\); the extract does not explicitly define it. The recurrence's index domain and boundary values for negative second indices are also unspecified.

For \(P(q)=\sum_{i\geq0}a_iq^i\), define its even and odd coefficient polynomials by
\[
P_{\mathrm{even}}(q)=\sum_{j\geq0}a_{2j}q^j,
\qquad
P_{\mathrm{odd}}(q)=\sum_{j\geq0}a_{2j+1}q^j.
\]
Call \(P\) \emph{parity log-concave} when both of these polynomials are log-concave. The extract states that the type \(B\) polynomials are not log-concave in general and gives the related Conjecture 4, attributed to [3, Conjecture 7.5]: every \(S_q^B[n,k]\) is parity log-concave for its admissible indices. This auxiliary conjecture retains the unresolved domain and boundary qualifications above.
\subsection{Short English statement}
The main question asks whether the coefficients of every source-defined \(q\)-Stirling polynomial \(S_q[n,k]\) form a log-concave sequence. The supplied record omits the definition and index domain of that family, so a fully specified conjecture requires further source clarification. The extract also records a type \(B\) parity log-concavity conjecture, computational evidence, and a reported asymptotic normality theorem.
\subsection{Sources}
\begin{itemize}
\item Supplied local UnsolvedMath registry record, permanent ID \(30005328\), source problem code \texttt{catalogue entry}, title \emph{Log-Concavity of q-Stirling Polynomials}, in the Oberwolfach Reports Open Problems collection. The record gives proposed year 2022 and supplies no named proposer.
\item Supplied report citation: \emph{Enumerative Combinatorics} (2023), Oberwolfach Reports, pp.\ 3215--3305, DOI \texttt{10.4171/owr/2022/57}. \url{https://doi.org/10.4171/owr/2022/57}. The background describes this collection as drawn from workshop reports and open-problem discussions.
\item Supplied literature source: \emph{q-Stirling numbers in type B}. \url{https://arxiv.org/abs/2205.14078}. The record does not explicitly identify this entry with the extract's numbered reference [3].
\item Internal attributions retained from the extract: [3, Conjecture 7.4], [3, Theorem 7.7], [3, Conjecture 7.5], Lorentzian polynomials [7], theory of atlases [6], and ideas in [10]. Their complete numbered bibliography is not supplied.
\item Literature assessment checked at \texttt{2026-08-22}: supplied dated curation, not independently verified. It labels the problem open and says that the stated log-concavity and asymptotic-normality questions were presented as conjectural, with no proof of the extracted claims found by that cutoff. This wording is in tension with the extract's report of an asymptotic normality theorem. The supplied curation requests independent verification of time-sensitive status and cited claims.
\end{itemize}
\end{document}
